% Source: ACM_Templates.pdf, PDF pp. 8-9, Manacher.
% Validation and changes: docs/validation.md; scripts/check_strings.py.
\subsection{Manacher：回文半径}
\label{algo:manacher}
\AlgorithmInfo
    {求最长回文子串、统计回文子串数量，并快速判断指定子串是否为回文。}
    {预处理时间、空间 $O(n)$；区间判断 $O(1)$，统计所有回文子串 $O(n)$。}
    {最长串起点和查询区间从 1 开始；内部变换串及半径数组从 0 开始。}

\subsubsection*{原理}
在原字符之间和两端插入分隔符，把奇、偶长度回文统一为同一种中心扩展。维护最右回文区间 $[l,r]$，当前中心在区间内时，先利用对称中心的半径，再检查超出已知范围的部分，整体只需线性次扩展。

\InlineCode{p[i]} 是含中心的半径，因此原串对应的回文长度为 $p_i-1$；以该中心产生的非空回文数为 $\lfloor p_i/2\rfloor$。原串区间 $[l,r]$ 对应变换串中心 $l+r-1$，比较半径与 $r-l+2$ 即可判断是否回文。

\subsubsection*{使用教程}
\InlineCode{init(s)} 预处理后，\InlineCode{len,pos} 为最长回文的长度和从 1 开始的起点。\InlineCode{check(l,r)} 返回闭区间是否为回文；\InlineCode{count()} 返回非空回文子串的出现次数，重复位置分别计数。换字符串重新 \InlineCode{init}，查询不改变预处理结果。
\begin{lstlisting}[style=mtycpp,numbers=none]
string s="cabbad";
Manacher ma;
ma.init(s);
string sub=s.substr(ma.pos-1,ma.len);
bool ok=ma.check(2,5);
long long cnt=ma.count();
\end{lstlisting}

\needspace{7\baselineskip}
\subsubsection*{模板代码}
\CodeFile{code/strings/manacher.cpp}

\subsubsection*{需要注意}
分隔符使用整数 $0$，原字符按无符号字节值加一编码，因此输入包含井号等字符也不冲突。\InlineCode{count()} 包含不同位置的重复回文，结果使用 64 位整数；空串的最长长度为 $0$，起点保留为 $1$。
